\documentclass[10pt]{ctexart}
\usepackage[a4paper,margin=1in]{geometry}
\usepackage{amsmath,amssymb,booktabs}
\usepackage{enumitem}
\usepackage{microtype}
\usepackage{hyperref}
\setlist[itemize]{leftmargin=1.5em}
\title{010：StepPPO 训练时间与资源消耗诊断}
\author{}
\date{}
\begin{document}
\maketitle
\vspace{-1.2em}

\section*{范围}
本文定位 StepPPO-v1 WebShop run 相比 GiGPO 更慢的原因，并区分 compute、memory、validation 与策略行为导致的开销。比较对象是 GiGPO seed 2026 与当前 StepPPO-v1 seed 2026 日志。

\section*{日志来源}
\begin{itemize}
\item GiGPO：\path{logs/webshop_gigpo_qwen25_15b_4gpu_paper_align_simple_20260530_114442.log}，完整到 step 250。
\item StepPPO-v1：\path{logs/webshop_step_ppo_qwen25_15b_4gpu_paper_align_simple_20260604_102927.log}，当前可用到 step 187。
\item 解析脚本：\path{VISULIZATION/plot_step_ppo_time_resource_diagnostics.py}。
\item 统计输出：\path{VISULIZATION/data/step_ppo_time_resource_summary.txt}。
\item 诊断图：\path{VISULIZATION/figures/fig_step_ppo_time_resource_diagnostics.png}。
\end{itemize}

\section*{整体时间}
non-validation training rows 的均值为：
\begin{center}
\begin{tabular}{lrrr}
\toprule
Metric & GiGPO & StepPPO-v1 & Ratio \\
\midrule
\path{timing_s/step} & 132.220 & 209.435 & 1.584 \\
\path{timing_s/gen} & 110.111 & 167.502 & 1.521 \\
\path{timing_s/old_log_prob} & 3.539 & 8.814 & 2.491 \\
\path{timing_s/ref} & 3.330 & 4.580 & 1.376 \\
\path{timing_s/update_actor} & 14.551 & 27.615 & 1.898 \\
\path{perf/throughput} & 2252.975 & 1836.456 & 0.815 \\
\bottomrule
\end{tabular}
\end{center}
平均来看，StepPPO 非验证 step 约为 GiGPO 的 $1.58\times$。在后期 steps 151--187，StepPPO 非验证 step 约为 GiGPO 的 $2.02\times$。

\section*{主因：Rollout Generation 变慢}
最大绝对时间差来自 rollout generation。平均 generation time 每个 non-validation step 增加约 $57.4$ 秒，占总 step time gap 的大部分。后期 segment 中，generation time gap 达到约 $87.0$ 秒。

这主要是策略行为导致的：StepPPO 生成更长、更容易截断、格式更差的 response：
\begin{center}
\begin{tabular}{lrrr}
\toprule
Metric & GiGPO & StepPPO-v1 & Ratio \\
\midrule
\path{response_length/mean} & 135.890 & 155.042 & 1.141 \\
\path{response_length/clip_ratio} & 0.004 & 0.026 & 5.825 \\
\path{global_seqlen/mean} & 289217.624 & 385207.740 & 1.332 \\
\path{episode/valid_action_ratio} & 0.978 & 0.803 & 0.821 \\
\bottomrule
\end{tabular}
\end{center}
更长 response 和更多 clipping 会同时增加 vLLM generation、old log prob、ref log prob 和 actor update 的 sequence length 成本。

\section*{实现额外开销}
StepPPO 确实存在直接实现开销。\path{verl/workers/actor/step_ppo_actor.py} 的 old-log-probability 路径在 value head 启用时不仅计算 token log probability，还会计算 old state values。actor update 路径也会重新计算 current values 并加入 clipped value loss。这解释了 old-log-probability $2.49\times$ 与 actor update $1.90\times$ 的开销。

另一个配置差异是 microbatch：StepPPO 4GPU 脚本使用 \path{actor_rollout_ref.actor.ppo_micro_batch_size_per_gpu=4}，GiGPO 脚本使用 $8$。更小 microbatch 降低显存压力，但增加 microbatch 数量，因此也会拖慢 update actor。

\section*{Validation 成本}
validation rows 更贵，因为 testing 本身要执行环境交互。平均 validation-step timing 为
\[
\mathrm{GiGPO}: 395.122\mathrm{s},
\qquad
\mathrm{StepPPO}: 548.947\mathrm{s}.
\]
平均 \path{timing_s/testing} 分别为 $265.629$ 秒和 $337.878$ 秒。因此比较 wall-clock 时必须保持 \path{trainer.test_freq} 一致。debug run 可以降低 validation 频率，但不能和 full protocol 的 wall-clock 直接比较。

\section*{资源消耗}
当前不是 memory-bound。StepPPO 的平均 GPU allocated memory 和 CPU memory 反而低于 GiGPO：
\begin{center}
\begin{tabular}{lrr}
\toprule
Metric & GiGPO & StepPPO-v1 \\
\midrule
\path{perf/max_memory_allocated_gb} & 71.710 & 60.905 \\
\path{perf/max_memory_reserved_gb} & 109.578 & 105.339 \\
\path{perf/cpu_memory_used_gb} & 551.620 & 433.726 \\
\bottomrule
\end{tabular}
\end{center}
这与 StepPPO 使用更小 actor microbatch 一致。主要瓶颈是 token generation 和额外 forward compute，而不是显存容量。

\section*{Wall-clock 粗估}
按观测均值外推，完整 250-step StepPPO-v1 约需 $19$ 小时；GiGPO seed 2026 约 $13$ 小时。如果后期 generation quality 继续退化，StepPPO 实际成本可能更高。

\section*{建议}
\begin{enumerate}
\item 先稳定 policy 行为。最大开销来自更长、更 invalid、更易 clipping 的 response。
\item 测试 \path{detach_value_backbone=True} 和更小 \path{value_loss_coef}，同时针对 stability 和 runtime。
\item 在显存允许时测试 StepPPO 的 \path{ppo_micro_batch_size_per_gpu=8}。
\item 优化 old-value collection，尽量把 log-probability 和 value computation 合并，避免额外 value forward。
\item 只在 smoke/debug 中降低 validation frequency；正式 wall-clock 对比必须保持一致。
\end{enumerate}

\section*{结论}
StepPPO 更慢有两个原因。主因是训练不稳定导致 rollout 行为变差，从而增加 response length、clipping 和 sequence length。次因是 value prediction/value loss 与较小 microbatch 带来的实现开销。当前日志显示它不是显存瓶颈。

\end{document}
